\documentclass[10pt,twocolumn]{article}

% ---- packages (standard TeX Live, compila con pdflatex) ----
\usepackage[T1]{fontenc}
\usepackage[utf8]{inputenc}
\usepackage{mathptmx}            % font serif tipo Times (look accademico)
\usepackage[margin=2cm]{geometry}
\usepackage{amsmath,amssymb}
\usepackage{booktabs}
\usepackage{titlesec}
\usepackage{fancyhdr}
\usepackage{enumitem}
\usepackage{siunitx}

% ---- titoli di sezione in maiuscoletto, centrati (stile esempio) ----
\titleformat{\section}[block]{\centering\scshape\large}{\thesection.}{0.5em}{}
\titleformat{\subsection}[block]{\bfseries}{\thesubsection}{0.5em}{}
\setlength{\columnsep}{0.8cm}

% ---- header corrente ----
\pagestyle{fancy}
\fancyhf{}
\fancyhead[R]{\small\itshape Tennis match prediction with machine learning}
\renewcommand{\headrulewidth}{0pt}

\title{\bfseries Can a Machine Learning Model Beat the Tennis Betting Market?\\
\large An Empirical Study with Public Data}
\author{%
  Author Name\\ Liceo Scientifico, Italy
}
\date{}

\begin{document}
\twocolumn[%
  \begin{@twocolumnfalse}
    \maketitle
    \begin{center}\textbf{Abstract}\end{center}
    \noindent
    This work investigates whether a machine learning model, trained only on
    free public data, can systematically profit from betting on ATP tennis
    matches. We build a complete pipeline---dynamic ELO ratings, leakage-free
    feature engineering, calibrated gradient boosting, and closing odds used as a
    feature---and evaluate it with rigorous statistical methods (walk-forward
    validation and bootstrap confidence intervals). The model reaches
    \textbf{67.6\% accuracy} and is well calibrated, yet its return on investment
    over more than \textbf{8{,}000 out-of-sample bets} is approximately
    \textbf{$-1\%$}, statistically indistinguishable from the bookmaker's margin.
    We conclude that the ATP tennis market, priced on sharp closing odds, is
    \emph{efficient} with respect to the information contained in public data:
    a negative but informative result, consistent with the efficient-market
    hypothesis. We emphasise the crucial distinction between predictive
    \emph{accuracy} and economic \emph{edge}, two notions that are frequently
    conflated.
    \vspace{0.6em}
  \end{@twocolumnfalse}
]

% ---------------------------------------------------------------
\begin{center}\scshape\large Indices\end{center}
\begin{enumerate}[leftmargin=1.4em,itemsep=2pt]
  \item[0.] \textbf{Abstract:} concise summary of the problem, method and findings.
  \item[1.] \textbf{Introduction:} why ``predicting'' and ``beating the market'' differ.
  \item[2.] \textbf{Related work:} ELO, efficient markets, prior tennis models.
  \item[3.] \textbf{Methodology:} data, ratings, leakage-free features, model,
            calibration, and how the edge is evaluated.
  \item[4.] \textbf{Results:} accuracy, calibration, ROI, statistical significance.
  \item[5.] \textbf{Discussion:} market efficiency, the hit-rate trap, limitations.
  \item[6.] \textbf{Conclusion:} implications of the negative result.
  \item[7.] \textbf{References.}
\end{enumerate}

% ---------------------------------------------------------------
\section{Introduction}
Sports betting markets are a natural laboratory for the efficient-market
hypothesis: every match has a price (the odds) that, if the market is efficient,
already incorporates all publicly available information. A predictive model is
therefore useful for betting only if it knows something the price does not.

This raises a question often misunderstood by hobbyists: is a model that
\emph{predicts winners accurately} the same as a model that \emph{beats the
market}? We argue, and demonstrate empirically, that the two are different. A
model can be highly accurate yet unprofitable, because the bookmaker prices the
odds around the same probabilities the model estimates, and keeps a margin.

We study men's professional tennis (ATP), where the data of Jeff Sackmann are
freely available and clean, and where closing odds (the last prices before a
match) are published for the main tour. Our hypotheses are:
\begin{itemize}[leftmargin=1.2em,itemsep=1pt]
  \item $H_0$ (efficient market): expected ROI $\le 0$.
  \item $H_1$: there exists a subset of matches with ROI $> 0$.
\end{itemize}

% ---------------------------------------------------------------
\section{Related Work}
The ELO rating system, originally designed for chess, has been widely applied to
tennis (e.g.\ FiveThirtyEight, Tennis Abstract), often specialised by surface.
The efficient-market hypothesis of Fama predicts that liquid betting markets, and
especially sharp bookmakers such as Pinnacle, leave little exploitable signal.
Prior machine-learning studies on tennis report accuracies around 65--70\%,
in line with the bookmakers' own implied accuracy, suggesting a hard ceiling
imposed by the intrinsic randomness of a single match (court conditions,
momentum, minor injuries).

% ---------------------------------------------------------------
\section{Methodology}

\subsection{Strategy}
We model the binary outcome ``does player~1 win?'' and convert the model
probability into a betting decision via the \emph{value} criterion (Sec.~3.7).
We deliberately use bookmaker odds in two distinct roles: (i) the closing
implied probability as an input \emph{feature}, and (ii) the market against which
we measure profit. Both are pre-match quantities, so neither introduces leakage.

\subsection{Data}
Match results come from the public Sackmann \texttt{tennis\_atp} repository:
about \num{63000} matches from 2005 to 2026, with per-match serve statistics,
rankings, surface and round. Closing odds come from \textit{tennis-data.co.uk}
(Pinnacle and Bet365), about \num{33000} matches from 2013 onward, ATP main tour
only. All data are ordered strictly by date.

\subsection{Dynamic ELO}
Each player has a general and a surface-specific rating. The expected win
probability of player $A$ against $B$ is
\begin{equation}
  E_A = \frac{1}{1 + 10^{\,(R_B - R_A)/400}},
\end{equation}
and ratings update after each match by
\begin{equation}
  R_A' = R_A + K\,(S_A - E_A),
\end{equation}
where $S_A\in\{0,1\}$ is the outcome and $K$ decreases with the number of matches
played (young players move faster than veterans).

\subsection{Leakage-free feature engineering}
\textbf{Golden rule:} for every match, features are computed from the state of
the world \emph{before} the match; the state is updated only \emph{after}.
Reversing this order leaks the outcome and produces fake accuracy. Features
include: ELO and surface-ELO differences, recent form (last 10 matches),
head-to-head, fatigue (minutes in the last 14 days), rolling serve statistics
(first-serve in \%, points won on serve, ace and double-fault rates,
break points saved), rankings and ranking points. The assignment of winner to
``player~1'' or ``player~2'' is randomised deterministically, so that the model
cannot learn the trivial rule ``player~1 always wins''.

\subsection{Model and calibration}
We train a gradient-boosted tree ensemble (XGBoost) with a strictly
\emph{temporal} split: training on the oldest seasons, validation on the next,
and test on the most recent (held-out). Raw boosting probabilities are
over-confident, so we apply \emph{isotonic regression} calibration fitted on the
validation set. Calibration does not change accuracy but makes the estimated edge
trustworthy. Quality of probabilities is measured by log-loss and the Brier
score,
\begin{equation}
  \mathrm{Brier} = \frac{1}{N}\sum_{i=1}^{N}\big(p_i - y_i\big)^2 .
\end{equation}

\subsection{Closing odds as a feature}
We add the overround-removed implied probability of the closing price as a model
feature. Given decimal odds $o_1,o_2$, the implied probability of player~1 is
\begin{equation}
  \hat{p}_1 = \frac{1/o_1}{1/o_1 + 1/o_2}.
\end{equation}
It is assigned to player~1 according to the deterministic swap, never according
to the result, so no leakage is introduced.

\subsection{Evaluating the edge}
For a bet at decimal odds $o$ on an outcome the model gives probability $p$, we
define the \emph{edge}
\begin{equation}
  \text{edge} = p - \frac{1}{o},
\end{equation}
and we bet only when $\text{edge} > \tau$ for a threshold $\tau$. The break-even
hit-rate at odds $o$ is exactly $1/o$. Stakes follow fractional Kelly,
\begin{equation}
  f^\star = \frac{p\,o - 1}{o - 1},
\end{equation}
capped for safety. To avoid fooling ourselves we use \textbf{walk-forward
validation}: the model is retrained on each past window and tested on the next,
never-seen window, with a \emph{fixed} filter applied throughout. Finally, a
\textbf{bootstrap} (\num{10000} resamples) gives a 95\% confidence interval (CI)
for the ROI,
\begin{equation}
  \text{ROI} = \frac{\sum_i \text{profit}_i}{\sum_i \text{stake}_i}.
\end{equation}

% ---------------------------------------------------------------
\section{Results}

\subsection{Accuracy and calibration}
Adding richer features and the closing-odds feature steadily improves the model
over a pure ELO baseline (Table~\ref{tab:acc}). The closing-odds feature alone
contributes roughly $+2$ accuracy points. The model is well calibrated: when it
says ``90\%'', the player wins about 91\% of the time.

\begin{table}[t]
\centering
\caption{Test-set performance (held-out seasons).}
\label{tab:acc}
\begin{tabular}{lccc}
\toprule
Model & Accuracy & LogLoss & Brier \\
\midrule
ELO baseline              & 63.4\% & 0.630 & 0.221 \\
XGBoost (ELO + serve)     & 65.0\% & 0.617 & 0.215 \\
\;+ closing odds (calib.) & \textbf{67.6\%} & 0.599 & 0.207 \\
\bottomrule
\end{tabular}
\end{table}

\subsection{Economic edge}
Despite the good accuracy, no betting strategy is profitable. Table~\ref{tab:roi}
reports the ROI and its 95\% bootstrap CI for several strategies evaluated
out-of-sample.

\begin{table}[t]
\centering
\caption{Out-of-sample betting results (Pinnacle closing odds).}
\label{tab:roi}
\begin{tabular}{lccc}
\toprule
Strategy & Bets & ROI & 95\% CI \\
\midrule
edge $\ge 0.04$ (WF)        & 8262 & $-1.1\%$ & $[-2.8,+0.7]$ \\
``safe'' ($p\ge0.80$)       & 768  & $-4.1\%$ & $[-7.6,-0.6]$ \\
rank band 31--50            & 1416 & $-1.9\%$ & $[-6.3,+2.6]$ \\
Slam+Masters, Hard/Grass    & 2572 & $+1.0\%$ & $[-1.9,+4.0]$ \\
\bottomrule
\end{tabular}
\end{table}

\noindent No strategy has a 95\% CI entirely above zero. The ``safe'' bets
(high-confidence picks) are in fact a \emph{proven loss}: the hit-rate is 82\%
but the average odds are only 1.21, whose break-even hit-rate is
$1/1.21 = 82.6\%$. The high hit-rate is eaten by the bookmaker's margin. This is
the central lesson:
\begin{equation}
  \text{edge} = \underbrace{p}_{\text{hit-rate}} - \frac{1}{o},
\end{equation}
so filtering by hit-rate alone, while ignoring the odds, selects short-priced
favourites and loses.

\subsection{The garden of forking paths}
The ``ranking 31--50'' segment looked promising at $+1.8\%$ over 95 bets in a
single backtest. Re-evaluated honestly over \num{1416} out-of-sample bets it
regressed to $-1.9\%$. This is a textbook example of a spurious pattern that
appears on small samples and vanishes on the full picture, and it justifies the
use of walk-forward validation with a pre-committed filter.

% ---------------------------------------------------------------
\section{Discussion}

\subsection{Market efficiency}
The result is consistent with an efficient market: Pinnacle's closing line
(2--3\% margin) already incorporates the information available in public data.
The model matches the market's accuracy but cannot exceed it; the residual
$-1\%$ ROI is essentially the margin.

\subsection{Limitations}
We use only free, pre-match data: no live (in-play) odds, and no private
information such as last-minute injuries or conditions. Historical odds exist
only for the ATP main tour, not for Challenger events where markets may be less
efficient. In production one bets hours before the match, not on the (sharper)
closing line, so the realistic edge would be even lower.

% ---------------------------------------------------------------
\section{Conclusion}
Using only free public data it is possible to build an accurate, well-calibrated
tennis predictor (67.6\%), but \emph{not} a profitable pre-match betting system
on the ATP market: the ROI is statistically indistinguishable from the
bookmaker's margin. The market is efficient with respect to this information.
The negative result is scientifically informative: it quantifies \emph{where
there is no value}, it makes explicit the difference between prediction and
profit, and it illustrates the methodological pitfalls (data leakage,
cherry-picking) of na\"ive data mining. Promising directions for future work are
less efficient markets (in-play, WTA, games totals and handicaps), the
integration of non-public data, and point-by-point Markov modelling.

% ---------------------------------------------------------------
\section{References}
\begin{enumerate}[leftmargin=1.6em,itemsep=1pt]
  \item J.~Sackmann, \textit{Tennis match datasets}, GitHub repository
        \texttt{tennis\_atp}.
  \item \textit{tennis-data.co.uk} --- historical closing odds.
  \item E.~Fama, ``Efficient Capital Markets: A Review of Theory and Empirical
        Work,'' \textit{Journal of Finance}, 1970.
  \item A.~Elo, \textit{The Rating of Chessplayers, Past and Present}, 1978.
  \item T.~Chen, C.~Guestrin, ``XGBoost: A Scalable Tree Boosting System,''
        \textit{KDD}, 2016.
  \item J.~L.~Kelly, ``A New Interpretation of Information Rate,''
        \textit{Bell System Technical Journal}, 1956.
  \item B.~Efron, ``Bootstrap Methods,'' \textit{Annals of Statistics}, 1979.
  \item A.~Gelman, E.~Loken, ``The Garden of Forking Paths,'' 2013.
\end{enumerate}

\end{document}
